% titel: Blatt 9, A6: Massenwirkungsgesetz aufstellen
% skript: Mortimer Kap. 14 (Massenwirkungsgesetz)

% Teil I: Werkzeuge, kein Lösungsweg.

Das Massenwirkungsgesetz setzt Produkte und Edukte einer Gleichgewichts-
reaktion zueinander ins Verhältnis.

\begin{werkzeugkasten}[Massenwirkungsgesetz]
Für $aA + bB \rightleftharpoons cC + dD$ (alle gasförmig oder in Lösung):
\[
  K_c = \frac{[C]^c [D]^d}{[A]^a [B]^b}.
\]
Wichtig:
\begin{itemize}
  \item Die \emph{Koeffizienten} der Reaktionsgleichung werden zu den
        \emph{Exponenten}.
  \item $K_c$ ist bei gegebener Temperatur konstant --- unabhängig von
        Ausgangskonzentrationen.
  \item Nur die Temperatur ändert den Wert von $K_c$; Konzentrationen
        ändern nur die Lage des Gleichgewichts, nicht $K_c$ selbst.
\end{itemize}
\textit{Skript:} Mortimer Kap. 14 (Massenwirkungsgesetz).
\end{werkzeugkasten}

\begin{achtung}[Darauf achten]
\begin{itemize}
  \item Steht in einer Teilfrage ``Wie ändert sich $K_c$ wenn [A]
        verdoppelt wird?`` --- Antwort: $K_c$ bleibt gleich (nur die Lage
        verschiebt sich, solange $T$ fest).
  \item $K_c$ \emph{verändern} kann man also nur über die Temperatur.
  \item Reine Feststoffe/Flüssigkeiten fliessen nicht in den Quotienten ein.
\end{itemize}
\end{achtung}
